\chapter{2024年全国甲卷（文）}

\section{单选题}

\begin{problem}
若 \(z=\sqrt{2}\i\)，则 \(z\bar z=\)（\quad）
\choices
    {\(-2\)}
    {\(-\sqrt{2}\)}
    {\(\sqrt{2}\)}
    {\(2\)}
\end{problem}

\begin{problem}
已知集合 \(A=\{1,2,3,4,5,9\}\)，\(B=\{x\mid x+1\in A\}\)，则 \(A\cap B=\)（\quad）
\choices
    {\(\{1,2,3\}\)}
    {\(\{3,4,9\}\)}
    {\(\{1,2,3,4\}\)}
    {\(\{2,3,4,5\}\)}
\end{problem}

\begin{problem}
设向量 \(\bs a=(x+1,x)\)，\(\bs b=(x,2)\)，\(\bs c=(1,1)\)，若 \(\bs a-2\bs b\) 与 \(\bs c\) 共线，则 \(x=\)（\quad）
\choices
    {\(\dfrac52\)}
    {\(\dfrac12\)}
    {\(-\dfrac12\)}
    {\(-\dfrac52\)}
\end{problem}

\begin{problem}
某独唱比赛的决赛阶段共有甲，乙，丙，丁四人参加，每人出场一次，出场次序由随机抽签确定. 则丙不是第一个出场，且甲或乙最后出场的概率是（\quad）
\choices
    {\(\dfrac16\)}
    {\(\dfrac14\)}
    {\(\dfrac13\)}
    {\(\dfrac12\)}
\end{problem}

\begin{problem}
记 \(S_n\) 为等差数列 \(\{a_n\}\) 的前 \(n\) 项和. 已知 \(S_9=1\)，则 \(a_3+a_7=\)（\quad）
\choices
    {\(\dfrac19\)}
    {\(\dfrac29\)}
    {\(\dfrac13\)}
    {\(\dfrac23\)}
\end{problem}

\begin{problem}
已知双曲线的两个焦点分别为 \((0,4)\)，\((0,-4)\)，点 \((-6,4)\) 在该双曲线上，则该双曲线的离心率为（\quad）
\choices
    {\(4\)}
    {\(3\)}
    {\(2\)}
    {\(\sqrt2\)}
\end{problem}

\begin{problem}
设函数 \(f(x)=\dfrac{\e^x+2\sin x}{1+x^2}\)，则曲线 \(y=f(x)\) 在点 \((0,1)\) 处的切线与两坐标轴所围成的三角形的面积为（\quad）
\choices
    {\(\dfrac16\)}
    {\(\dfrac13\)}
    {\(\dfrac12\)}
    {\(\dfrac23\)}
\end{problem}

\begin{problem}
函数 \(y=-x^2+(\e^x-\e^{-x})\sin x\) 在区间 \([-2.8,2.8]\) 的图象大致为（\quad）
\choices
    {\choicebitmap[width=0.15\paperwidth]{img/2024/national_paper_a_liberal/q8_optA.png}}
    {\choicebitmap[width=0.15\paperwidth]{img/2024/national_paper_a_liberal/q8_optB.png}}
    {\choicebitmap[width=0.15\paperwidth]{img/2024/national_paper_a_liberal/q8_optC.png}}
    {\choicebitmap[width=0.15\paperwidth]{img/2024/national_paper_a_liberal/q8_optD.png}}
\end{problem}

\begin{problem}
直线 \(2x-y-2=0\) 与圆 \(x^2+y^2-6x-8y=0\) 交于 \(A\)，\(B\) 两点，则 \(|AB|=\)（\quad）
\choices
    {\(4\)}
    {\(5\)}
    {\(8\)}
    {\(10\)}
\end{problem}

\begin{problem}
已知 \(\dfrac{\cos\alpha}{\cos\alpha-\sin\alpha}=\sqrt3\)，则 \(\tan\left(\alpha+\dfrac\pi4\right)=\)（\quad）
\choices
    {\(2\sqrt3+1\)}
    {\(2\sqrt3-1\)}
    {\(\dfrac{\sqrt3}{2}\)}
    {\(1-\sqrt3\)}
\end{problem}

\begin{problem}
设 \(\alpha\)，\(\beta\) 为两个平面，\(m\)，\(n\) 为两条直线，且 \(\alpha\cap\beta=m\). 下述四个命题：
\begin{circlelist}
\item 若 \(m\parallel n\)，则 \(n\parallel\alpha\) 或 \(n\parallel\beta\)；
\item 若 \(m\perp n\)，则 \(n\perp\alpha\) 或 \(n\perp\beta\)；
\item 若 \(n\parallel\alpha\) 且 \(n\parallel\beta\)，则 \(m\parallel n\)；
\item 若 \(n\) 与 \(\alpha\)，\(\beta\) 所成的角相等，则 \(m\perp n\).
\end{circlelist}
其中所有真命题的编号是（\quad）
\choices
    {\(\circled{1}\circled{3}\)}
    {\(\circled{2}\circled{4}\)}
    {\(\circled{1}\circled{2}\circled{3}\)}
    {\(\circled{1}\circled{3}\circled{4}\)}
\end{problem}

\begin{problem}
记 \(\triangle ABC\) 的内角 \(A\)，\(B\)，\(C\) 的对边分别为 \(a\)，\(b\)，\(c\)，已知 \(B=60^\circ\)，\(b^2=\dfrac94ac\)，则 \(\sin A+\sin C=\)（\quad）
\choices
    {\(\dfrac32\)}
    {\(\sqrt2\)}
    {\(\dfrac{\sqrt7}{2}\)}
    {\(\dfrac{\sqrt3}{2}\)}
\end{problem}

\section{填空题}

\begin{problem}
函数 \(f(x)=\sin x-\sqrt3\cos x\) 在区间 \([0,\pi]\) 的最大值为\fillinblank{}.
\end{problem}

\begin{problem}
已知圆台甲，乙的上底面半径均为 \(r_1\)，下底面半径均为 \(r_2\)，圆台甲，乙的母线长分别为 \(2(r_2-r_1)\)，\(3(r_2-r_1)\)，则圆台甲与乙的体积之比为\fillinblank{}.
\end{problem}

\begin{problem}
已知 \(a>1\) 且 \(\dfrac1{\log_8a}-\dfrac1{\log_a4}=-\dfrac52\)，则 \(a=\)\fillinblank{}.
\end{problem}

\begin{problem}
当 \(x>0\) 时，曲线 \(y=x^3-3x\) 与曲线 \(y=-(x-1)^2+a\) 有两个交点，则 \(a\) 的取值范围是\fillinblank{}.
\end{problem}

\section{解答题}

\begin{problem}
记 \(S_n\) 为等比数列 \(\{a_n\}\) 的前 \(n\) 项和，已知 \(2S_n=3a_{n+1}-3\).
\begin{enumerate}
\item 求 \(\{a_n\}\) 的通项公式；
\item 求数列 \(\{S_n\}\) 的前 \(n\) 项和.
\end{enumerate}
\end{problem}

\begin{problem}
某工厂进行生产线智能化升级改造. 升级改造后，从该工厂甲，乙两个车间的产品中随机抽取 \(150\text{ 件}\) 进行检验，数据如下：
\begin{center}
\begin{tabular}{c|ccc|c}
 & 优级品 & 合格品 & 不合格品 & 总计\\ \hline
甲车间 & 26 & 24 & 0 & 50\\
乙车间 & 70 & 28 & 2 & 100\\ \hline
总计 & 96 & 52 & 2 & 150
\end{tabular}
\end{center}
\begin{enumerate}
\item 填写如下列联表：
\begin{center}
\begin{tabular}{c|cc}
 & 优级品 & 非优级品\\ \hline
甲车间 &\fillinblank{} &\fillinblank{} \\
乙车间 &\fillinblank{} &\fillinblank{} \\
\end{tabular}
\end{center}
能否有 \(95\%\) 的把握认为甲，乙两车间产品的优级品率存在差异？能否有 \(99\%\) 的把握认为甲，乙两车间产品的优级品率存在差异？
\item 已知升级改造前该工厂产品的优级品率 \(p=0.5\). 设 \(\bar p\) 为升级改造后抽取的 \(n\text{ 件}\) 产品的优级品率，如果 \(\bar p>p+1.65\sqrt{\dfrac{p(1-p)}n}\)，则认为该工厂产品的优级品率提高了. 根据抽取的 \(150\text{ 件}\) 产品的数据，能否认为生产线智能化升级改造后，该工厂产品的优级品率提高了？（\(\sqrt{150}\approx12.247\)）
\end{enumerate}
附：\(K^2=\dfrac{n(ad-bc)^2}{(a+b)(c+d)(a+c)(b+d)}\)，\(P(K^2\ge k)\) 的临界值如下：
\begin{center}
\begin{tabular}{c|ccc}
 & 0.050 & 0.010 & 0.001\\ \hline
\(k\) & 3.841 & 6.635 & 10.828
\end{tabular}
\end{center}
\end{problem}

\begin{problem}
如图，在以 \(A\)，\(B\)，\(C\)，\(D\)，\(E\)，\(F\) 为顶点的五面体中，四边形 \(ABCD\) 与四边形 \(ADEF\) 均为等腰梯形，\(EF\parallel AD\)，\(BC\parallel AD\)，\(AD=4\)，\(AB=BC=EF=2\)，\(ED=\sqrt{10}\)，\(FB=2\sqrt3\)，\(M\) 为 \(AD\) 的中点.
\begin{enumerate}
\item 证明：\(BM\parallel\) 平面 \(CDE\)；
\item 求 \(M\) 到平面 \(FAB\) 的距离.
\end{enumerate}

\bitmapfigure[max width=0.70\linewidth]{img/2024/national_paper_a_liberal/q19_fig1.png}
\end{problem}

\begin{problem}
已知函数 \(f(x)=a(x-1)-\ln x+1\).
\begin{enumerate}
\item 讨论 \(f(x)\) 的单调性；
\item 设 \(a\le2\)，证明：当 \(x>1\) 时，\(f(x)<\e^{x-1}\).
\end{enumerate}
\end{problem}

\begin{problem}
设椭圆 \(C:\dfrac{x^2}{a^2}+\dfrac{y^2}{b^2}=1\)（\(a>b>0\)） 的右焦点为 \(F\)，点 \(M\left(1,\dfrac32\right)\) 在 \(C\) 上，且 \(MF\perp x\) 轴.
\begin{enumerate}
\item 求 \(C\) 的方程；
\item 过点 \(P(4,0)\) 的直线交 \(C\) 于 \(A\)，\(B\) 两点，\(N\) 为线段 \(FP\) 的中点，直线 \(NB\) 交直线 \(MF\) 于点 \(Q\)，证明：\(AQ\perp y\) 轴.
\end{enumerate}
\end{problem}

\section{选考题：解答题}

\begin{problem}
在平面直角坐标系 \(xOy\) 中，以坐标原点 \(O\) 为极点，\(x\) 轴正半轴为极轴建立极坐标系，曲线 \(C\) 的极坐标方程为 \(\rho=\rho\cos\theta+1\).
\begin{enumerate}
\item 写出 \(C\) 的直角坐标方程；
\item 设直线 \(l\) 的参数方程为 \(\begin{cases}
x=t,\\
y=t+a
\end{cases}\ (t\text{ 为参数})\)，若 \(C\) 与 \(l\) 相交于 \(A\)，\(B\) 两点，且 \(|AB|=2\)，求 \(a\).
\end{enumerate}
\end{problem}

\begin{problem}
已知实数 \(a\)，\(b\) 满足 \(a+b\ge3\).
\begin{enumerate}
\item 证明：\(2a^2+2b^2>a+b\)；
\item 证明：\(|a-2b^2|+|b-2a^2|\ge6\).
\end{enumerate}
\end{problem}

